% chapters/07/contoh.tex
\section{Dua Contoh Analisis}
\label{sec:contoh-bab-7}

\subsection{Contoh Pertama: Menghitung Benda yang Menembus Bongkahan}

\begin{contohbox}
Sebuah benda bergerak dengan kecepatan 900 piksel per detik pada laju 30
bingkai per detik. Perpindahannya pada setiap bingkai adalah

\[
900 \div 30 = 30 \text{ piksel}.
\]

Bongkahan yang dilaluinya tebalnya 20 piksel. Perhatikan apa yang terjadi pada
kedua bingkai yang berurutan.
\end{contohbox}

\begin{table}[htbp]
\centering
\caption{Posisi benda terhadap bongkahan setebal 20 piksel pada dua bingkai berurutan}
\label{tab:tunneling}
{\setlength{\tabcolsep}{4pt}\footnotesize%
\begin{tabularx}{\textwidth}{@{}C{1.6cm} C{3.4cm} C{3.6cm} Y@{}}
\toprule
\headrow \textbf{Bingkai} & \textbf{Sisi kanan benda} & \textbf{Kedudukan terhadap bongkahan} & \textbf{Hasil pemeriksaan} \\
\midrule
Sebelum & 415 piksel & 15 piksel sebelum tepi bongkahan pada 430 & Tidak beririsan. \\
\addlinespace
Berikutnya & 445 piksel & 15 piksel sesudah tepi jauh bongkahan pada 450 & Tidak beririsan. \\
\bottomrule
\end{tabularx}}
\tabsource{Bongkahan menempati 430 sampai 450 piksel. Pada kedua bingkai yang
diperiksa, benda berada di luar bongkahan sehingga tidak pernah terdeteksi
bersentuhan.}
\end{table}

Benda itu tidak pernah beririsan dengan bongkahan pada saat pemeriksaan
dilakukan, walaupun lintasannya jelas melintasi bongkahan. Perpindahan 30
piksel pada satu bingkai lebih panjang daripada tebal bongkahan 20 piksel,
sehingga benda ``melompati'' bongkahan itu.

Tiga cara perbaikan berikut dapat dipilih, dan ketiganya sah.

\begin{enumerate}
  \item \textbf{Memperkecil panjang langkah.} Membatasi selang waktu maksimum
        yang boleh dipakai pada satu bingkai, misalnya pada 1/60 detik.
        Perpindahan pada contoh di atas turun menjadi 15 piksel, dan bongkahan
        setebal 20 piksel kembali terdeteksi.
  \item \textbf{Membagi perpindahan menjadi beberapa langkah kecil.} Satu
        bingkai yang panjang dipecah menjadi beberapa langkah yang
        masing-masing pendek. Cara ini lebih tepat daripada cara pertama,
        tetapi memakan lebih banyak perhitungan.
  \item \textbf{Mempertebal bongkahan.} Cara paling sederhana, dan cukup
        untuk permainan yang kecepatannya terbatas. Bongkahan setebal 40
        piksel tidak lagi dapat dilalui benda pada contoh di atas.
\end{enumerate}

\subsection{Contoh Kedua: Analisis Tabrakan pada Sisi yang Salah}

\begin{contohbox}
Seorang pemain berjalan ke kanan dan menabrak sisi kiri sebuah bongkahan.
Pemain berukuran 30 kali 42 piksel, dan bongkahan menempati 430 sampai 550
piksel mendatar serta 310 sampai 330 piksel tegak.

Pada bingkai yang diperiksa, \code{rect} pemain bernilai (404, 300, 30, 42),
sehingga menempati 404 sampai 434 piksel mendatar dan 300 sampai 342 piksel
tegak. Kedua persegi beririsan, sebab mendatarnya bertumpang tindih pada 430
sampai 434, dan tegaknya pada 310 sampai 330.
\end{contohbox}

\begin{table}[htbp]
\centering
\caption{Akibat urutan pemeriksaan terhadap letak pemain}
\label{tab:sumbu-salah}
{\setlength{\tabcolsep}{4pt}\footnotesize%
\begin{tabularx}{\textwidth}{@{}L{4.4cm} C{2.6cm} C{2.6cm} Y@{}}
\toprule
\headrow \textbf{Cara pemeriksaan} & \textbf{Sisi yang diselesaikan} & \textbf{Titik tengah pemain} & \textbf{Akibat yang terlihat} \\
\midrule
Kedua sumbu sekaligus, sumbu tegak didahulukan &
Atas bongkahan & 321 menjadi 289 piksel &
Pemain terangkat 32 piksel ke atas bongkahan, walaupun datang dari samping. \\
\addlinespace
Sumbu mendatar lebih dahulu &
Sisi kiri bongkahan & 415 menjadi 415 piksel &
Pemain berhenti tepat di sisi bongkahan, sebagaimana mestinya. \\
\bottomrule
\end{tabularx}}
\tabsource{Posisi mendatar titik tengah pemain semula 415 piksel, dan posisi
tegaknya 321 piksel.}
\end{table}

Pada baris pertama, program memilih penyelesaian pada sumbu
tegak karena kecepatan tegak pemain bernilai positif: pemain memang sedang
jatuh. Pilihan itu masuk akal bagi program, tetapi salah bagi pemain yang
datang dari samping.

Gejala inilah yang paling sering dikeluhkan pemula sebagai ``pemain tiba-tiba
terbang ke atas bongkahan''. Sekali diketahui penyebabnya, perbaikannya
sederhana: periksa sumbu mendatar lebih dahulu, selesaikan sampai tuntas, baru
periksa sumbu tegak. Dengan urutan itu, pada saat sumbu tegak diperiksa, pemain
sudah tidak lagi bertumpang tindih dengan bongkahan, sehingga penyelesaian pada
sumbu tegak hanya terjadi ketika pemain memang datang dari atas.
